\section{Problem Statement}

Let \texttt{V} be the vulnerabilities observable as of a decision time \texttt{t0} and \texttt{H} the
hosts in a fleet. The candidate set of \textbf{vulnerability-host pairs} is
\texttt{P = {(v, h) : v in V, h in H, v applies to h}}, with applicability determined by
software/configuration matching. Each pair carries exploit-intelligence, severity,
asset-criticality, and local-exposure attributes observable at \texttt{t0} under strict
no-future-leakage (Section 7). A \textbf{strategy} induces a total order over \texttt{P}. A
maintenance window admits at most a fixed \textbf{capacity} \texttt{c} pair-actions, subject to
blackout windows, operational constraints, and approval gates. The \textbf{objective} is
to evaluate and compare strategies under this capacity constraint using operational
and ranking metrics (Section 9), recording each decision as tamper-evident
evidence.

\textbf{Non-goals.} We do not perform autonomous patching (the scheduler is a
capacity-constrained scheduling simulation modeling approval and timing, not patch
execution or success); we do not perform production validation (the fleet is
public-sector-shaped synthetic); we do not benchmark commercial RBVM products; and
we do not claim the framework guarantees or proves compliance --- it produces
evidence that supports review.
